\begin{examplebox}
	\textbf{\textcolor{CExample}{例 1（基础）}}
	\textbf{\textcolor{CExample}{题目：}}
	已知向量 $\vec{a}=(3,4)$，求 $|\vec{a}|$。
	\textbf{\textcolor{CExample}{解题步骤：}}
	\begin{description}[style=nextline, leftmargin=2.8em, labelsep=0.5em, itemsep=0.45em]
		\item[\textbf{第一步（代入公式）：}]
			直接运用向量模长公式：
			\[
				|\vec{a}| = \sqrt{3^2+4^2}.
			\]
			\par\textbf{理由：} 已知 $x=3,y=4$，代入公式。
		\item[\textbf{第二步（计算平方和）：}]
			计算根号内的平方和：
			\[
				\begin{aligned}
					3^2+4^2 &= 9+16 \\ &= 25.
			\end{aligned}
			\]
			\par\textbf{理由：} 平方运算与加法。
		\item[\textbf{第三步（开方求值）：}]
			开平方得到模长：
			\[
				|\vec{a}| = \sqrt{25}.
			\]
			\par\textbf{理由：} $\sqrt{25}=5$，模长取正值，故 $|\vec{a}|=5$。
	\end{description}
	\textbf{\textcolor{CExample}{关键结论：}}
	\textbf{$|\vec{a}|=5$}。

	\medskip
	\textbf{\textcolor{CExample}{例 2（稍变形）}}
	\textbf{\textcolor{CExample}{题目：}}
	已知两点 $A(1,2),B(4,6)$，求 $|\overrightarrow{AB}|$。
	\textbf{\textcolor{CExample}{解题步骤：}}
	\begin{description}[style=nextline, leftmargin=2.8em, labelsep=0.5em, itemsep=0.45em]
		\item[\textbf{第一步（坐标差向量）：}]
			写出向量 $\overrightarrow{AB}$ 的坐标并化简：
			\[
				\begin{aligned}
					\overrightarrow{AB} &= (4-1,6-2) \\ &= (3,4).
			\end{aligned}
			\]
			\par\textbf{理由：} 向量坐标等于终点减起点。
		\item[\textbf{第二步（代入公式）：}]
			代入模长公式：
			\[
				|\overrightarrow{AB}| = \sqrt{3^2+4^2}.
			\]
			\par\textbf{理由：} 由模长公式直接写出。
		\item[\textbf{第三步（化简计算）：}]
			化简求值：
			\[
				\begin{aligned}
					|\overrightarrow{AB}| &= \sqrt{9+16} \\ &= \sqrt{25} \\ &= 5.
			\end{aligned}
			\]
			\par\textbf{理由：} 先计算平方和，再开方。
	\end{description}
	\textbf{\textcolor{CExample}{关键结论：}}
	\textbf{$|\overrightarrow{AB}|=5$}。

	\medskip
	\textbf{\textcolor{CExample}{例 3（综合应用）}}
	\textbf{\textcolor{CExample}{题目：}}
	已知点 $A(1,k),B(4,2)$，且 $|\overrightarrow{AB}|=5$，求 $k$ 的值。
	\textbf{\textcolor{CExample}{解题步骤：}}
	\begin{description}[style=nextline, leftmargin=2.8em, labelsep=0.5em, itemsep=0.45em]
		\item[\textbf{第一步（坐标差向量）：}]
			写出 $\overrightarrow{AB}$ 的坐标：
			\[
				\begin{aligned}
					\overrightarrow{AB} &= (4-1,2-k) \\ &= (3,2-k).
			\end{aligned}
			\]
			\par\textbf{理由：} 终点坐标减起点坐标。
		\item[\textbf{第二步（列方程）：}]
			根据模长公式列方程：
			\[
				\sqrt{3^2 + (2-k)^2} = 5.
			\]
			\par\textbf{理由：} 代入模长公式，模长等于5。
		\item[\textbf{第三步（求解参数）：}]
			解方程求 $k$：
			\[
				\begin{aligned}
					\sqrt{9 + (2-k)^2} &= 5 \\ 9 + (2-k)^2 &= 25 \\ (2-k)^2 &= 16 \\ 2-k &= \pm 4 \\ k &= 2 \mp 4.
			\end{aligned}
			\]
			\par\textbf{理由：} 两边平方，移项，开平方。
	\end{description}
	\textbf{\textcolor{CExample}{关键结论：}}
	\textbf{$k=6$ 或 $k=-2$}。
\end{examplebox}
